% Generated draft: requires mathematical and source review.
\documentclass[11pt]{article}
\usepackage[margin=25mm]{geometry}
\usepackage{fontspec}
\setmainfont{Latin Modern Roman}
\usepackage{amsmath,amssymb,mathtools,unicode-math}
\setmathfont{Latin Modern Math}
\usepackage{enumitem,needspace,xurl,hyperref}
\setlength{\parindent}{0pt}
\setlength{\parskip}{5pt}
\setlength{\emergencystretch}{3em}
\newcommand{\N}{\mathbb{N}}
\newcommand{\Z}{\mathbb{Z}}
\newcommand{\Q}{\mathbb{Q}}
\newcommand{\R}{\mathbb{R}}
\newcommand{\C}{\mathbb{C}}
\newcommand{\F}{\mathbb{F}}
\newcommand{\A}{\mathbb{A}}
\newcommand{\PP}{\mathbb{P}}
\newcommand{\E}{\mathbb{E}}
\newcommand{\eps}{\varepsilon}
\newcommand{\dd}{\,\mathrm{d}}
\begin{document}
\section*{UnsolvedMath ID: 10400071}
\textbf{Source problem code:} catalogue entry
\subsection{Definitions and mathematical statement}
Reading the flattened exponents as reciprocals, let $A(t)$ satisfy $A(1)=1$ and $A(t)=A(t^{-1})$. The natural nontrivial interpretation is a Laurent polynomial, and set
\[
R_A=\mathbb Q[t^{\pm1},A(t)^{-1}],\qquad R_0=\mathbb Q[t^{\pm1}].
\]
The first ring consists of Laurent-polynomial numerators divided by nonnegative powers of $A(t)$. If ``polynomial'' means only nonnegative powers and the reciprocal symmetry is literal, it forces $A=1$; the intended Laurent/normalization convention therefore matters.

For each loop degree $l$, find an explicit $\mathbb Q$-basis of
\[
\mathcal A_{R_A}(\varnothing;\mathbb Q)^{(\mathrm{loop}\ l)},
\]
and, in particular, of
\[
\mathcal A_{R_0}(\varnothing;\mathbb Q)^{(\mathrm{loop}\ l)}.
\]
The source intends a space of labeled Jacobi diagrams with empty skeleton. In a usual beaded-graph presentation, edges carry ring labels and vertices carry cyclic-order data; orientation reversal changes labels through the involution $t\mapsto t^{-1}$. However, the full label, AS/IHX, and any vertex/gauge relations are not supplied. Nor does the extract specify connected versus disconnected diagrams or the precise loop grading; first Betti number is the usual context.

The requested bases are over $\mathbb Q$, not assumed free-module bases over $R_A$. An explicit spanning and independent family is sought, not merely an abstract assertion that a vector-space basis exists. The general localization and unlocalized case remain separate tasks until their exact diagram conventions are fixed.
\subsection{Short English statement}
Find explicit rational bases in every loop degree for closed Jacobi-diagram spaces labeled by $\mathbb Q[t^{\pm1},1/A(t)]$, especially $\mathbb Q[t^{\pm1}]$. Clarify reciprocal/Laurent polynomial notation and the diagram relations.
\subsection{Sources}
\begin{itemize}
\item Ohtsuki, \emph{Problems on invariants of knots and 3-manifolds} (2002), Problem 3.21, PDF page 70, \url{https://msp.org/gtm/2002/04/gtm-2002-04-024s.pdf}; accessed 2026-07-29, PDF-font-structure extraction.
\item Local registry \texttt{data/unsolvedmath/problems.json}, permanent UnsolvedMath ID $10400071$, source code catalogue entry, \emph{Problem 3.21}. Dataset attribution: Ulam AI and the UnsolvedMath Contributors.
\end{itemize}
\end{document}
